\documentclass[10pt]{amsart}

\usepackage[margin=1in]{geometry}
\usepackage[T1]{fontenc}
\usepackage{lmodern}
\usepackage{amsmath,amssymb,amsthm,mathtools}
\usepackage{microtype}
\usepackage{booktabs}
\usepackage{longtable}
\usepackage{url}
\usepackage{array}
\usepackage{enumitem}
\usepackage[dvipsnames]{xcolor}
\usepackage[colorlinks=true,linkcolor=MidnightBlue,citecolor=OliveGreen,urlcolor=MidnightBlue]{hyperref}
\usepackage[capitalize,nameinlink]{cleveref}

\setlist{leftmargin=*,itemsep=2pt,topsep=4pt}
\emergencystretch=2.5em
\setlength{\tabcolsep}{4pt}
\renewcommand{\arraystretch}{1.12}

% ---------------------------------------------------------------- theorems
\theoremstyle{plain}
\newtheorem{theorem}{Theorem}[section]
\newtheorem{lemma}[theorem]{Lemma}
\newtheorem{proposition}[theorem]{Proposition}
\newtheorem{corollary}[theorem]{Corollary}
\newtheorem{importedlemma}[theorem]{Imported Lemma}
\newtheorem{importedtheorem}[theorem]{Imported Theorem}
\newtheorem*{hypA}{Hypothesis $(\mathrm{H}_{\mathrm{A}})$}
\newtheorem*{hypO}{Hypothesis $(\mathrm{H}_{\mathrm{O}})$}
\theoremstyle{definition}
\newtheorem{definition}[theorem]{Definition}
\newtheorem{conjecture}[theorem]{Conjecture}
\theoremstyle{remark}
\newtheorem{remark}[theorem]{Remark}

\crefname{importedlemma}{Imported Lemma}{Imported Lemmas}
\crefname{importedtheorem}{Imported Theorem}{Imported Theorems}
\crefname{conjecture}{Conjecture}{Conjectures}
\numberwithin{equation}{section}

% ---------------------------------------------------------------- macros
\newcommand{\Z}{\mathbb{Z}}
\newcommand{\Q}{\mathbb{Q}}
\newcommand{\Zw}{\Z[\omega]}
\newcommand{\poly}{\operatorname{poly}}
\newcommand{\Kron}{\operatorname{Kron}}
\newcommand{\gap}{\operatorname{gap}}
\newcommand{\cG}{\mathcal{G}}
\newcommand{\cB}{\mathcal{B}}
\newcommand{\cL}{\mathcal{L}}
\newcommand{\Ind}[1]{[\,#1\,]}
\newcommand{\ket}[1]{|#1\rangle}
\newcommand{\bra}[1]{\langle#1|}
\newcommand{\amp}[3]{\langle #1|#2|#3\rangle}
\newcommand{\A}[1][]{\textup{[A\if\relax\detokenize{#1}\relax\else, #1\fi]}}
\newcommand{\Oa}[1][]{\textup{[O\if\relax\detokenize{#1}\relax\else, #1\fi]}}

\begin{document}

\title[Exact quantum circuit amplitudes after the ETH refutations]{Exact amplitudes of sparse quantum circuits\\ in subexponential time,\\ assuming the new $\#$SAT algorithms}

\author{Vaibhav Tulsyan}

\date{October 2026}

\subjclass[2020]{68Q12, 68Q25, 81P68}
\keywords{Quantum circuit simulation, strong simulation, Exponential Time Hypothesis, fine-grained complexity, quantum supremacy, tensor networks, Clifford+T}

\begin{abstract}
In October 2026 OpenAI and Anthropic claimed deterministic subexponential-time algorithms for
$\#k$-SAT, refuting the Exponential Time Hypothesis. We record a corollary for quantum circuits.
Assuming three internal results of the Anthropic manuscript, the amplitude $\amp{y}{C}{x}$ of a
Clifford+T circuit with $G$ two-qubit gates in arbitrary connectivity and $G_1$ single-qubit gates
can be computed exactly, as an element of $\Z[\omega]/\sqrt2^{\,k}$ with $\omega=e^{i\pi/4}$, in
deterministic time $2^{O(G/\sqrt{\log G})}\poly(n+G_1)$; the same holds for exact marginals and
hence exact sampling, and for depth $d$ the bound is $2^{O(dn/\sqrt{\log n})}$. Through OpenAI's
tensor-contraction theorem the bounds hold in polynomial space when the number of Hadamard-type
gates is $O(G)$. The proof is a sum over paths with one bit per wire segment, read as an
Anthropic-type network, with a Kronecker substitution for $\Z[\omega]$ and a Galois bound on the
coefficients. Conversely, constant-depth circuits with $O(n)$ gates encode $\#3$-SAT with $O(n)$
clauses, so the rate matches that of the refutations. We check the fine-grained
quantum-supremacy conjectures of Morimae--Tamaki, Dalzell--Harrow--Koh--La Placa,
Huang--Newman--Szegedy, Aaronson--Chen and Buhrman--Patro--Speelman against their sources and
record which are refuted, which lose their premise, and which survive; a counting form of the
Calabro--Impagliazzo--Paturi width reduction refutes the gap conjecture of Morimae and Tamaki's
depth-scaling paper. Dense statements (cubic phases, permanents, random circuits) survive, and we
explain why. Nothing here concerns
approximate sampling, the constants are galactic, and everything is conditional on the claimed
refutations, which are under review.
\end{abstract}

\maketitle


% =====================================================================
\section{Introduction}\label{sec:intro}
% =====================================================================

\subsection{Background}
The Exponential Time Hypothesis (ETH) of Impagliazzo and Paturi~\cite{ImpagliazzoPaturi01}
asserts that there is $\delta>0$ such that $3$-SAT on $n$ variables cannot be solved in time
$O(2^{\delta n})$; by the sparsification lemma~\cite{IPZ01} it is equivalent to the same statement
for the number $m$ of clauses.
In October 2026 two manuscripts appeared, by OpenAI~\cite{OpenAI26} and by
Anthropic~\cite{Anthropic26}, each claiming a deterministic algorithm that counts the satisfying
assignments of a $k$-CNF in subexponential time, and hence a refutation of ETH. Both come with
Lean~4 formalisations of their top-level statements. These are \emph{claimed} results, still
under review by the community. We write \A{} and \Oa{} for the two manuscripts. The Main Theorem
of \A{} counts the models of a $k$-CNF with $m$ clauses on $n$ variables deterministically in time
$2^{O(\min\{km/\sqrt{\log m},\;k^2n\log\log n/\sqrt{\log n}\})}$; the main theorem of \Oa{} does
so in $(n+m)^{O_k(1)}2^{O_k(n\log\log n/\sqrt{\log n})}$ bit operations and polynomial space, and
in $2^{O_{k,D}(n/\sqrt{\log n})}$ when $m\le Dn$.

A part of the case for ``quantum computational supremacy'' rests on fine-grained hypotheses of
exactly this kind. Morimae and Tamaki~\cite{MT19a,MT19b,MT20}, Huang, Newman and
Szegedy~\cite{HNS18,HNS19} and Dalzell, Harrow, Koh and La Placa~\cite{DHKL20} derive lower
bounds on the classical simulation of quantum circuits from ETH, SETH and relatives, usually by
writing $\#\mathrm{SAT}$ or the gap $\gap(f)=\sum_x(-1)^{f(x)}$ of a Boolean function as a circuit
amplitude. The question we ask is the converse: what do the new algorithms give for the classical
simulation of quantum circuits, and which of these lower bounds survive?

\subsection{Results}
A circuit $C$ acts on $n$ qubits and is a sequence of gates from the set
\[
  \cG=\{H,T,T^\dagger,S,S^\dagger,X,Z\}\cup\{\mathrm{CNOT},\mathrm{CZ},\mathrm{CS},\mathrm{SWAP},\mathrm{CH}\}.
\]
It has $G$ two-qubit gates, in arbitrary connectivity, and $G_1$ single-qubit gates. Put
$\omega=e^{i\pi/4}$, so that $\Zw=\Z[x]/(x^4+1)$, $i=\omega^2$ and $\sqrt2=\omega-\omega^3\in\Zw$.
Every gate $U\in\cG$ is stored as a pair $(N_U,\kappa_U)$ with $N_U$ a matrix over $\Zw$ and
$U=N_U/\sqrt2^{\,\kappa_U}$, where $\kappa_U=1$ for $H$ and $\mathrm{CH}$ and $\kappa_U=0$
otherwise (for $\mathrm{CH}$, $N=\mathrm{diag}(\sqrt2,\sqrt2)\oplus\left(\begin{smallmatrix}1&1\\1&-1\end{smallmatrix}\right)$). Put
$k(C):=\sum_U\kappa_U$, the number of $H$ and $\mathrm{CH}$ gates. Running times are counted as in
\A: operations on words long enough to hold an address, plus table queries; all logarithms are to
base~$2$. The input has length $\poly(n+G+G_1)$.

\begin{hypA}
The following statements of \A{} hold as stated there: Lemma~8, case~(i) (the reduction from
general networks to cliques, \cref{imp:lemma8}); Theorem~A(a),(b) (one algorithm for the
$(K_\ell,\lambda)$-cliques of all levels $\ell$, \cref{imp:thmA}); and Theorem~B (pieces with
parameters $(1/\sqrt2,c)$, \cref{imp:thmB}).
\end{hypA}

\begin{hypO}
The Boolean tensor-contraction theorem in the appendix of \Oa{} holds as stated there
(\cref{imp:tensor}).
\end{hypO}

\begin{theorem}[Exact amplitudes; machine-checked]\label{thm:main}
Assume $(\mathrm{H}_{\mathrm{A}})$. There are an absolute constant $c$ and one deterministic
Turing machine that, given a circuit $C$ over $\cG$ on $n$ qubits with $G$ two-qubit gates, and
$x,y\in\{0,1\}^n$, encoded as a word of length $L$ with $\log_2(L+2)\le(G+2)^{0.9}$, outputs
integers $a_0,a_1,a_2,a_3$ with
\[
  \amp{y}{C}{x}=\frac{a_0+a_1\omega+a_2\omega^2+a_3\omega^3}{\sqrt2^{\,k(C)}}
\]
within $\lfloor c\cdot2^{c(G+2)/\sqrt{\log_2(G+2)}}\rfloor$ steps. Without the condition on $L$,
the same holds within $\lfloor c\cdot2^{cG/\sqrt{\log_2(G+2)}}(L+2)^{c(G+2)^{1/15}}\rfloor$ steps,
and for every $\delta>0$ there is a machine running within $\lfloor c_\delta\,2^{\delta G}(L+2)^{c_\delta}\rfloor$ steps.
\end{theorem}
\noindent The condition on $L$ holds, for example, whenever $n(G+G_1)\le 2^{G^{0.9}}$; on paper,
the proof below also gives $2^{cG/\sqrt{\log G}}\poly(n+G_1)$ under the weaker condition
$\log(G_1+2)\le G^{0.9}$, but that variant is not machine-checked and not claimed.

\begin{theorem}[Polynomial space; written proof, not machine-checked]\label{thm:openai}
Assume $(\mathrm{H}_{\mathrm{O}})$. Put $s:=\max\{4G,\,k(C),\,4\}$. The output of
\cref{thm:main} can be computed in $\poly(n+G+G_1)\,2^{O(s/\sqrt{\log s})}$ bit operations and
$\poly(n+G+G_1)$ bits of space. In particular, if $k(C)=O(G)$ (for example, if the number of
$H$ and $\mathrm{CH}$ gates is $O(G)$), the time is $2^{O(G/\sqrt{\log G})}\poly(n+G_1)$ and the
space is polynomial.
\end{theorem}

\begin{corollary}[Marginals and exact sampling; written proof, not machine-checked]\label{cor:marginals}
Assume $(\mathrm{H}_{\mathrm{A}})$, and let $C,x$ be as in \cref{thm:main}. For every
$S\subseteq[n]$ and $y_S\in\{0,1\}^S$, the marginal probability
$p(y_S)=\sum_{y_{[n]\setminus S}}|\amp{y}{C}{x}|^2$ can be computed exactly, as
$(r_0+r_1\sqrt2)/2^{k(C)}$ with $r_0,r_1\in\Z$, within $2^{c'G/\sqrt{\log G}}\poly(n+G_1)$
operations, for an absolute constant $c'$. Consequently there is a Las Vegas algorithm that
outputs $y\in\{0,1\}^n$ with probability exactly $|\amp{y}{C}{x}|^2$ in expected time
$n\cdot2^{c'G/\sqrt{\log G}}\poly(n+G_1)$. Under $(\mathrm{H}_{\mathrm{O}})$ the same holds in
polynomial space when $k(C)=O(G)$.
\end{corollary}

\begin{corollary}[Constant and slowly growing depth]\label{cor:depth}
Assume $(\mathrm{H}_{\mathrm{A}})$. Amplitudes, marginals and exact samples of depth-$d$ circuits
over $\cG$ on $n\ge4$ qubits can be computed within $2^{O(dn/\sqrt{\log n})}$ operations; the
exponent is $o(n)$ whenever $d=o(\sqrt{\log n})$.
\end{corollary}

\begin{proposition}[The rate is that of $\#3$-SAT]\label{prop:tight}
There is a polynomial-time map from $3$-CNFs $\varphi$ with $n$ variables and $m$ clauses to
circuits $C_\varphi$ over $\cG$ of depth $O(1)$ with $O(n+m)$ qubits and $O(n+m)$ gates, and to
an integer $n'$, such that $\amp{0}{C_\varphi}{0}=\#\varphi/2^{n'}$. Hence an algorithm for exact
amplitudes of $G$-gate circuits in time $T(G)$ counts the models of $3$-CNFs with $m\ge n$
clauses in time $T(O(m))\poly(m)$.
\end{proposition}

So, assuming the refutations, the exact amplitude problem for circuits with $O(n)$ gates is
exactly as hard as $\#3$-SAT with $O(n)$ clauses, up to constants in the exponent: a bound
$2^{o(G/\sqrt{\log G})}$ in \cref{thm:main} would improve the rate of \A{} for $m=O(n)$. Before
2026 the same reduction showed, under the counting version $\#$ETH~\cite{DHMTW14}, that
$2^{\Omega(n)}$ time was needed. \Cref{tab:sota} compares \cref{thm:main} with the classical
methods for exact strong simulation.

\begin{table}[ht]
\centering
\small
\caption{Exact strong simulation of a circuit with $n$ qubits, $G$ two-qubit gates and depth
$d$, arbitrary connectivity. The last two rows are conditional.}\label{tab:sota}
\begin{tabular}{@{}>{\raggedright\arraybackslash}p{4.4cm}>{\raggedright\arraybackslash}p{4.6cm}>{\raggedright\arraybackslash}p{2.0cm}>{\raggedright\arraybackslash}p{3.6cm}@{}}
\toprule
Method & Time & Space & Source\\
\midrule
state vector & $2^n\poly$ & $2^n$ & folklore\\
Feynman path sum & $2^{O(G)}$ & $\poly$ & folklore\\
recursive (Savitch-type) & $O(n(2d)^{n+1})$ & $O(n\log d)$ & \cite[Thm.~4.1]{AC17}\\
tensor contraction & $2^{O(\mathrm{tw})}$, $\mathrm{tw}=\Theta(n)$ possible at $d=O(1)$ & $2^{O(\mathrm{tw})}$ & \cite{MarkovShi08}\\
geometrically local, dim.\ $\ge2$ & subexponential in $G$ & -- & \cite{WahlStrelchuk}\\
\cref{thm:main} & $2^{O(G/\sqrt{\log G})}\le2^{O(dn/\sqrt{\log n})}$ & as time & under $(\mathrm{H}_{\mathrm{A}})$\\
\cref{thm:openai} & same, if $k(C)=O(G)$ & $\poly$ & under $(\mathrm{H}_{\mathrm{O}})$\\
\bottomrule
\end{tabular}
\end{table}

For depth $d=O(1)$ in arbitrary connectivity we know of no previous $2^{o(n)}$ exact method:
the interaction graph of such a circuit can be a bounded-degree expander, whose treewidth is
$\Theta(n)$. Against the polynomial-space bound $(2d)^{O(n)}=2^{O(n\log d)}$ of Aaronson and
Chen~\cite{AC17}, \cref{thm:openai} gives $2^{O(dn/\sqrt{\log n})}$ in the worst case $G=\Theta(dn)$,
which is asymptotically smaller only when $d/\log d=O(\sqrt{\log n})$.

\subsection{Contrast with BosonSampling}
BosonSampling amplitudes are permanents of $n\times n$ submatrices of Haar-random
unitaries~\cite{AaronsonArkhipov13}, and Ryser's $2^n\poly(n)$ formula~\cite{Ryser63} is still
the best exact method. Nothing above applies to them. A qubit circuit is a sum over \emph{binary}
variables, one per wire segment, coupled by bounded-arity local tables, so it has $O(G)$ shared
bits. A dense permanent is a sum over \emph{permutations}; written as an interferometer it has
$\Theta(m^2)$ beam splitters on modes carrying up to $n$ photons, and any local encoding seems to
need about $\log n$ bits per row. Informal notes in the project repository, not claimed here, show that
``row-local'' networks for the permanent in the model of \A{} (rows grouped into clusters of at
most $s$) have total at least $n\log(n/s)$, too much for the engine to beat $2^n$; only sparse
matrices go through. The conjecture per-int-NSETH of~\cite{DHKL20} is therefore
untouched (\cref{sec:table}).

\subsection{What this does not do}
\begin{itemize}
\item \emph{Approximate sampling; many gates.} Random-circuit sampling rests on the non-collapse
of the polynomial hierarchy and on average-case conjectures about approximate amplitudes (QUATH,
HOG~\cite{AC17}) for circuits with $G\ge n^2$ gates; none of this is touched. For
$G\ge n\sqrt{\log n}$ the bound is no better than $2^n$.
\item \emph{Practical simulation.} The constants are galactic. Using \A[Theorem~A(a)] with the
smallest admissible $x_0=18$, the leading term of the exponent is $e_\ell U$ with $U\approx4G$
when most wire segments cross the partition, and
$15/\sqrt{\log G}\le e_\ell\le114/\sqrt{\log G}$ at the level chosen in \cref{sec:time}. For
depth $d$, so $G\approx dn/2$, the leading term alone is below $n$ only if $\log n$ exceeds
roughly $900\,d^2$ (with the lower estimate) to $5.2\cdot10^4\,d^2$ (with the upper one), and
lower-order terms only make this worse. \Oa{} gives no explicit constants. Experiments with
$50$ to $1000$ qubits are not affected.
\end{itemize}

\subsection{What is new here, and what is not}
Neither manuscript mentions quantum circuits: in the source of \Oa{} the words ``quantum'',
``qubit'' and ``amplitude'' do not occur, and \A{} mentions tensor networks, including Markov and
Shi's paper on simulating quantum computation by tensor contraction~\cite{MarkovShi08}, only in
its related work. Nor did we find the corollary in the press or discussion threads we searched.
Mathematically, however, \cref{thm:openai} is an immediate corollary of
the appendix theorem of \Oa, which contracts Boolean tensor networks with one-bit wires: anyone
who reads that theorem next to Markov and Shi will see it. What we add is the care needed to make
it a theorem: the coefficient ring ($\Zw$ is not a subring of $\Z$, nor of the truncated $2$-adic
ring of \Oa), the precision needed, which comes from a Galois-conjugate bound
(\cref{lem:galois}), the resulting restriction $k(C)=O(G)$ in the route through \Oa{} and its
absence in the route through \A, marginals and exact sampling, the converse reduction, the
check of the fine-grained supremacy literature against the primary sources
(\cref{sec:table}), including a counting form of the width reduction of Calabro, Impagliazzo and
Paturi (\cref{lem:width}), and the analysis of cubic phases (\cref{sec:cubic}). We claim the
careful statement and proof, not the idea.

% =====================================================================
\section{The imported results}\label{sec:imported}
% =====================================================================

We recall the definitions and results of \A{} that we use, in its notation except where a letter
clashes with ours, as in~\cite{Tulsyan26}. Conventions of \A: $\log=\log_2$ and
$[x=u]\in\{0,1\}$.

\begin{definition}[{\A[Definition~1]}]\label{def:factor}
A function $F\colon X\times X'\to\Z$ \emph{factors through $c$ bits} if there are
$f\colon X\times\{0,1\}^c\to\Z$ and $f'\colon X'\times\{0,1\}^c\to\Z$ with
$F(z,z')=\sum_{\theta\in\{0,1\}^c} f(z,\theta)f'(z',\theta)$. In particular, for any
$h\colon X\to\{0,1\}^c$ and $h'\colon X'\to\{0,1\}^c$ the function
$[h(z)=h'(z')]=\sum_\theta [h(z)=\theta][h'(z')=\theta]$ factors through $c$ bits.
\end{definition}

\begin{definition}[{\A[Definition~2]}]\label{def:network}
Let $P,B,\tau\ge1$ be integers. A \emph{network on $J$ with parameters $P,B,\tau$} consists of a
finite set $J$ of nodes; for each node $a$ a finite set $S_a$ of settings and a function
$W_a\colon S_a\to\Z$; and for each unordered pair $\{a,b\}$ of distinct nodes a function
$\Psi_{a,b}\colon S_a\times S_b\to\Z$ (with $\Psi_{b,a}(z_b,z_a)=\Psi_{a,b}(z_a,z_b)$), an integer
$c_{a,b}=c_{b,a}\ge0$ and a factorisation $(f_{a,b},f_{b,a})$ of $\Psi_{a,b}$ through $c_{a,b}$
bits. The settings of $a$ are encoded as bit strings of one fixed length $P_a\le P$, and
\begin{enumerate}[label=\textup{(\roman*)}]
\item $\sum_{b\ne a} c_{a,b}\le P$ for every $a\in J$;
\item all values of all $W_a$, $\Psi_{a,b}$, $f_{a,b}$ are below $2^B$ in absolute value;
\item an algorithm receives $P$, $B$ and the numbers $P_a$, $c_{a,b}$ as input, and can decide
membership in $S_a$ and compute any value of $W_a$, $\Psi_{a,b}$, $f_{a,b}$ with at most $\tau$
operations.
\end{enumerate}
The \emph{value} of the network is
$Z=\sum_{z\in\prod_a S_a}\prod_{a\in J}W_a(z_a)\prod_{\{a,b\}}\Psi_{a,b}(z_a,z_b)$, and its
\emph{total} is $U=\sum_{(a,b)\colon a\ne b}c_{a,b}$.
\end{definition}

In the \emph{$(K,\lambda)$-clique} of \A[Definition~3] each of $K$ nodes holds a $\lambda$-bit
string for each other node, $c_{w,w'}=\lambda$, $\Psi_{w,w'}=[z_w(w')=z_{w'}(w)]$, and the tables
$|W_w|<2^B$ are queried at unit cost.


\begin{importedlemma}[{\A[Lemma~8, case~(i)]}]\label{imp:lemma8}
Let $K\ge2$, $g=3K^8$ and $\zeta=1/(2K)$. Take a network on $[g]$ with parameters $P,B,\tau$ and
total $U$, and put $L=\lceil(U+gP/K)/K^2\rceil$ and $B'=\lceil 2gP/K\rceil+g^2B$.
If an algorithm outputs the value of every $(K,L)$-clique with $|W_w|<2^{B'}$ within $T$
operations, then the value of the network is computed within
\begin{equation}\label{eq:lemma8}
  \bigl(K^g+T\cdot 2^{2gP/K}\cdot\tau\bigr)\cdot\bigl(g(P+B)\bigr)^{C_1}
\end{equation}
operations, by one procedure, for an absolute constant $C_1$.
\end{importedlemma}

\begin{importedtheorem}[{\A[Theorem~A]}]\label{imp:thmA}
Suppose there are pieces with parameters $(\rho,c)$ in the sense of \A[Definition~7], and let
$\alpha'=\log(1/\rho)$. Fix an integer $x_0\ge18$ with $2^{x_0}\ge c+5$ and define
$K_0=2^{x_0}$, $\zeta_\ell=2^{-\ell-5}$, $g_\ell=2^{3\ell+x_0}K_\ell$, $K_{\ell+1}=g_\ell^2$,
$e_0=\tfrac12$ and $e_{\ell+1}=(\rho+3\zeta_\ell)e_\ell+3/K_\ell$.
\begin{enumerate}[label=\textup{(\alph*)}]
\item $\log K_\ell=(3x_0+6)2^\ell-6\ell-2x_0-6$, so $K_\ell<K_{\ell+1}\le K_\ell^{10}$; and
$\frac12\rho^\ell\le e_\ell\le 0.85\,\rho^\ell$, where $\rho^\ell\le((3x_0+6)/\log K_\ell)^{\alpha'}$.
\item One deterministic algorithm, given $\ell\ge0$, $\lambda,B\ge1$ and query access to tables
with $|W_w|<2^B$, outputs the value of the $(K_\ell,\lambda)$-clique within
\begin{equation}\label{eq:thmA}
  2^{e_\ell K_\ell^2\lambda}\cdot\bigl(K_\ell(\lambda+B)\bigr)^{C_0K_\ell^2}
\end{equation}
operations, with $C_0$ depending only on $c$.
\end{enumerate}
\end{importedtheorem}

\begin{importedtheorem}[{\A[Theorem~B]}]\label{imp:thmB}
There are pieces with parameters $(1/\sqrt2,c)$ for an absolute constant $c$.
\end{importedtheorem}

\begin{importedtheorem}[{\Oa[Appendix, Boolean tensor contraction], up to notation}]\label{imp:tensor}
A Boolean tensor network is a bipartite multigraph in which each edge carries a single bit; each
vertex $v$ has an ordered list $\delta(v)$ of incident edges and a tensor $T_v$ assigning a ring
element to each Boolean assignment to those edges; its contraction is
$Z=\sum_{x\in\{0,1\}^{E}}\prod_vT_v(x|_{\delta(v)})$, an isolated vertex contributing its scalar
entry. Fix positive constants $C,A,b$. Let $s\ge2$ be an integer, and let $\cL$ be the total
length of the graph and tensor descriptions of a bipartite tensor network with at most $Cs$
edges and maximum degree at most $\lceil s^{9/10}\rceil$. Suppose that the number of vertices is
bounded by a fixed polynomial in $s$, and that the entries belong to
\[
  \cB_{P_0,\ell}=(\Z/2^{P_0}\Z)[V]/(V^\ell),\qquad 1\le P_0,\ell\le Cs.
\]
Suppose, uniformly for every vertex and Boolean choice of its ports, that its entry can be
computed from the description in time $(s+\cL)^b2^{As^{9/10}}$ and space $(s+\cL)^b$. Then the
contraction can be computed deterministically in $(s+\cL)^{O(1)}2^{O(s/\sqrt{\log s})}$ bit
operations and $(s+\cL)^{O(1)}$ bits of space, with constants depending only on $C,A,b$.
\end{importedtheorem}

\A's Lean development proves \cref{imp:lemma8} (in a form weaker by a polynomial factor),
\cref{imp:thmA}(b) and \cref{imp:thmB} for arbitrary networks in its own format, and our Lean proof
of \cref{thm:main} applies them directly to the circuit network (\cref{sec:formalisation}). The
appendix of \Oa{} is not formalised, so \cref{thm:openai} is proved on paper only.

% =====================================================================
\section{Proof of \texorpdfstring{\cref{thm:main}}{Theorem 1.1}}\label{sec:proof}
% =====================================================================

\subsection{The sum over paths on wire segments}\label{sec:paths}
For a qubit $q$ let $\gamma_{q,1},\dots,\gamma_{q,r_q}$ be the two-qubit gates acting on $q$, in
circuit order; $q$ is \emph{touched} if $r_q\ge1$. Every single-qubit gate on a touched line is
multiplied into the next two-qubit gate on that line, and the single-qubit gates after
$\gamma_{q,r_q}$ are multiplied into $\gamma_{q,r_q}$ from the left. This is legitimate because a
single-qubit gate on line $q$ commutes with every gate that does not act on $q$. The result is a
list of $G$ \emph{merged} gates $\hat\Gamma_i=\hat N_i/\sqrt2^{\,\hat\kappa_i}$ with
$\hat N_i\in M_4(\Zw)$, where $\hat\kappa_i$ is the sum of the $\kappa$ of the gates merged into
it. An untouched line contributes the scalar $\amp{y_q}{R_q}{x_q}=u_q/\sqrt2^{\,\kappa_q}$, where
$R_q$ is the product of its gates. The merging costs $\poly(n+G+G_1)$ operations, since all
coefficients have $O(k(C)+G+G_1)$ bits.

A \emph{segment} is a piece of line $q$ between $\gamma_{q,j}$ and $\gamma_{q,j+1}$; it carries
one bit $\sigma_e$. The input leg of $\gamma_{q,1}$ is fixed to $x_q$ and the output leg of
$\gamma_{q,r_q}$ to $y_q$. With $n_t$ touched lines there are $s_0=\sum_q(r_q-1)=2G-n_t\le2G$
segments. For a merged gate $i$ write $\sigma_{\mathrm{in}}(i),\sigma_{\mathrm{out}}(i)\in\{0,1\}^2$
for the values on its input and output legs (a segment bit or a fixed $x_q$, $y_q$).

\begin{lemma}[Sum over paths]\label{lem:paths}
Put $k_t:=\sum_i\hat\kappa_i$ and $A:=\sum_{\sigma\in\{0,1\}^{s_0}}\prod_{i=1}^G\hat N_i[\sigma_{\mathrm{out}}(i),\sigma_{\mathrm{in}}(i)]\in\Zw$. Then
\[
 \amp{y}{C}{x}=\Bigl(\prod_{q\text{ untouched}}u_q\Bigr)\cdot\frac{A}{\sqrt2^{\,k(C)}},\qquad
 k(C)=k_t+\sum_{q\text{ untouched}}\kappa_q,
\]
and $A=\sqrt2^{\,k_t}\amp{y_t}{C_t}{x_t}$, where $C_t$ is the circuit $C$ restricted to the
touched lines.
\end{lemma}

\begin{proof}
Insert $\sum_{b\in\{0,1\}}\ket b\bra b$ on every segment and expand the matrix product.
\end{proof}

\subsection{Coefficients: a Galois bound and a Kronecker map}
For $\alpha=\sum_ja_j\omega^j\in\Zw$ write $\|\alpha\|_\infty=\max_j|a_j|$ and
$\|\alpha\|_1=\sum_j|a_j|$. Multiplication by $\omega^j$ permutes the coefficients up to sign, so
$\|\alpha\beta\|_1\le\|\alpha\|_1\|\beta\|_1$.

\begin{lemma}[Galois bound]\label{lem:galois}
Let $M$ be a product of gates from $\cG$ (on any number of qubits), let $\kappa$ be the sum of
their $\kappa_U$, and let $N$ be the product of their numerators, so $M=N/\sqrt2^{\,\kappa}$. Then
every entry $\alpha$ of $N$ satisfies $\|\alpha\|_\infty\le2^{\kappa/2}$.
\end{lemma}

\begin{proof}
The Galois group of $\Q(\omega)/\Q$ consists of $\sigma_t\colon\omega\mapsto\omega^t$, $t\in\{1,3,5,7\}$;
it is abelian and contains complex conjugation $\sigma_7$. Hence, for a unitary $U$ with entries
in $\Q(\omega)$, $\sigma_t(U)\sigma_t(U)^*=\sigma_t(UU^*)=I$, so $\sigma_t(U)$ is unitary. Applying
$\sigma_t$ entrywise to all factors, and using $\sigma_t(\sqrt2)=\pm\sqrt2$, we get
$\sigma_t(N)=(\pm\sqrt2)^{\kappa}\sigma_t(M)$ with $\sigma_t(M)$ unitary, so
$|\sigma_t(\alpha)|\le2^{\kappa/2}$. Finally $\sum_{t}\omega^{rt}=4[r\equiv0]$ for
$r\in\{-3,\dots,3\}$, which gives $a_j=\frac14\sum_t\sigma_t(\alpha)\,\omega^{-jt}$ and
$|a_j|\le2^{\kappa/2}$.
\end{proof}

By \cref{lem:paths} and \cref{lem:galois} (applied to $C_t$), $\|A\|_\infty\le2^{k_t/2}$, and every
entry of every $\hat N_i$ has $\|\cdot\|_\infty\le2^{\hat\kappa_i/2}$, hence $\|\cdot\|_1\le4\cdot2^{\hat\kappa_i/2}$.

\begin{lemma}[Kronecker map]\label{lem:kron}
Let $b\ge2$ and $M_b:=2^{4b}+1$. The substitution $x\mapsto2^b$ defines a ring homomorphism
$\kappa_b\colon\Zw=\Z[x]/(x^4+1)\to\Z/M_b\Z$, because $2^{4b}\equiv-1\pmod{M_b}$. It is injective
on $\{\alpha:\|\alpha\|_\infty\le2^{b-1}-1\}$, and $\alpha$ is recovered from $\kappa_b(\alpha)$ by
taking the representative in $(-M_b/2,M_b/2]$ and reading its balanced base-$2^b$ digits.
\end{lemma}

\begin{proof}
For $\alpha$ in the box, $|\sum_ja_j2^{bj}|\le(2^{b-1}-1)(1+2^b+2^{2b}+2^{3b})<2^{4b-1}<M_b/2$, and
balanced digits are unique.
\end{proof}

Put $\Kron_b(\alpha):=\sum_ja_j2^{bj}\in\Z$, an integer representative of $\kappa_b(\alpha)$, and
fix $b:=\lceil k_t/2\rceil+2$, so that $\|A\|_\infty\le2^{b-1}-1$.

No ring embedding is available ($x^4=-1$ has no solution modulo $4$, nor modulo $V$); a
homomorphism injective on a box containing the answer suffices.

\subsection{The network}\label{sec:network}
Fix $g\ge1$ and a partition of the merged gates into groups $V_1,\dots,V_g$ with
$|V_a|\le\lceil G/g\rceil$ (some groups may be empty). A segment is \emph{internal} to $a$ if both
its end gates lie in $V_a$, and \emph{crossing} otherwise. Let $\mathrm{int}_a$ be the number of
segments internal to $a$ and $P_a$ the number of crossing segments with an end in $V_a$.

\begin{definition}\label{def:qnet}
The network $\mathcal N(C,x,y,(V_a),b)$ on $[g]$ has settings $S_a=\{0,1\}^{P_a}$, the bits of the
crossing segments at $a$;
\[
 W_a(z_a):=\Kron_b\Bigl(\sum_{\text{internal bits of }a}\ \prod_{i\in V_a}\hat N_i[\sigma_{\mathrm{out}}(i),\sigma_{\mathrm{in}}(i)]\Bigr),
\]
the sum taken in $\Zw$ over the $2^{\mathrm{int}_a}$ values of the internal segment bits, with the
crossing bits read from $z_a$; and
$\Psi_{a,b}(z_a,z_b)=\prod_{e}[z_a(e)=z_b(e)]$ over the segments $e$ crossing between $a$ and $b$,
which factors through $c_{a,b}:=\#\{e\text{ crossing between }a\text{ and }b\}$ bits by
\cref{def:factor}, with $f_{a,b}\in\{0,1\}$.
\end{definition}

\begin{proposition}[Parameters and value]\label{prop:params}
Let $P:=4\lceil G/g\rceil+1$, $B:=3b+2+4\lceil G/g\rceil+\lceil k_t/2\rceil$ and
$\tau:=2^{2\lceil G/g\rceil}\poly(G,k(C))$. Then $\mathcal N=\mathcal N(C,x,y,(V_a),b)$ is a
network on $[g]$ with parameters $P,B,\tau$ in the sense of \cref{def:network}, its total is
$U\le 4G$, and its value $Z$ satisfies $Z\equiv\kappa_b(A)\pmod{M_b}$. Hence $A$ is obtained from
$Z$ by \cref{lem:kron}.
\end{proposition}

\begin{proof}
\emph{Condition (i).} $\sum_{b\ne a}c_{a,b}=P_a$, and every merged gate has four legs, so
$P_a\le4|V_a|\le P$. This is where bounded gate arity enters; nothing about connectivity is used.
\emph{The total.} $U=\sum_aP_a=2\cdot\#\{\text{crossing segments}\}\le2s_0\le4G$.
\emph{Condition (ii).} The values of $\Psi$ and $f$ are $0$ or $1$. For $W_a$, the inner sum
$\alpha_a$ has $\|\alpha_a\|_\infty\le\|\alpha_a\|_1\le2^{\mathrm{int}_a}\prod_{i\in V_a}4\cdot2^{\hat\kappa_i/2}$,
and $\mathrm{int}_a\le2|V_a|$, so $\log\|\alpha_a\|_\infty\le4\lceil G/g\rceil+k_t/2$; and
$|\Kron_b(\alpha)|<\|\alpha\|_\infty2^{3b+1}$. So $|W_a|<2^B$.
\emph{Condition (iii).} Membership in $S_a$ is trivial, $\Psi$ and $f$ compare bits, and $W_a(z_a)$
is a sum of $2^{\mathrm{int}_a}\le2^{2\lceil G/g\rceil}$ products of $|V_a|$ elements of $\Zw$ with
$O(k(C)+G)$-bit coefficients, read from the precomputed merged gates.
\emph{Value.} $\Kron_b(\alpha)\equiv\kappa_b(\alpha)\pmod{M_b}$ and $\kappa_b$ is a ring
homomorphism, so
$Z\equiv\kappa_b\bigl(\sum_z\prod_a\alpha_a(z_a)\prod_{\{a,b\}}\Psi_{a,b}\bigr)$. The factors
$\Psi$ restrict the sum to settings that agree on every crossing segment, and the inner sums run
over the internal segments, so the argument of $\kappa_b$ is $A$.
\end{proof}

\subsection{Running time}\label{sec:time}
The algorithm merges the gates, computes $k_t$, $b$ and the untouched scalars, chooses a level
$\ell$ (below), puts $K:=K_\ell$ and $g:=3K^8$, takes any partition into groups of size at most
$\lceil G/g\rceil$, computes $P,B,P_a,c_{a,b}$, and runs the procedure of \cref{imp:lemma8} on
$\mathcal N$ with the algorithm of \cref{imp:thmA}(b) at level $\ell$ as the clique algorithm.
It reduces the value $Z$ modulo $M_b$, decodes $A$ by \cref{lem:kron}, and multiplies by the
untouched scalars.

\emph{Hypotheses.} \Cref{imp:lemma8} needs a network on $[g]$ with integer parameters
$P,B,\tau\ge1$ and $g=3K^8$; \cref{prop:params} supplies it. It then needs an algorithm for every
$(K,L)$-clique with $|W|<2^{B'}$: \cref{imp:thmA}(b) at level $\ell$, with $\lambda=L\ge1$ and
$B=B'\ge1$, gives $T=2^{e_\ell K^2L}(K(L+B'))^{C_0K^2}$. \Cref{imp:thmA} needs pieces, which
\cref{imp:thmB} supplies with $\rho=1/\sqrt2$, $\alpha'=\frac12$, and $c$, $x_0$ absolute. The
circuit reaches the clique algorithm only through the tables that \cref{imp:lemma8} builds from
the $W_a$ and $f_{a,b}$, and \cref{imp:lemma8} itself searches the $K^g$ groupings of the $g$ nodes
for one with small pairwise totals, so no balancing beyond the partition is needed.

\emph{The exponent.} By \eqref{eq:lemma8} and \eqref{eq:thmA}, with $e:=e_\ell\le1$ and
$K^2L\le U+gP/K+K^2$, the cost is $2^X$ with
\begin{equation}\label{eq:exponent}
  X\le1+\max\Bigl\{g\log K,\ e\bigl(U+\tfrac{gP}{K}+K^2\bigr)+C_0K^2\log\bigl(K(L+B')\bigr)+\tfrac{2gP}{K}+\log\tau\Bigr\}
  +C_1\log\bigl(g(P+B)\bigr).
\end{equation}
For $G<G_0:=K_0^{30}$, an absolute constant, the algorithm instead sums over the $\le2G$ segment
bits directly, in $O(1)\cdot\poly(n+G_1)$ operations. For $G\ge G_0$ let $\ell$ be the largest
index with $K_\ell^{30}\le G$. Then $K_{\ell+1}>G^{1/30}$, and $K_{\ell+1}\le K_\ell^{10}$ gives
$G^{1/300}<K\le G^{1/30}$, so $g=3K^8\le3G^{8/30}$ and, by \cref{imp:thmA}(a),
\[
  e\le0.85\Bigl(\frac{3x_0+6}{\log K}\Bigr)^{1/2}\le\frac{0.85\,(300(3x_0+6))^{1/2}}{\sqrt{\log G}}=:\frac{c_e}{\sqrt{\log G}}.
\]
Also $gP\le g(4G/g+5)=4G+5g$, $L\le U+gP+1=O(G)$ and
$B'\le2gP+1+g^2B=O(G^{1.6}(1+k(C)))$. Since $k(C)\le G+G_1$ and $\log(G_1+2)\le G^{0.9}$, we have
$\log(k(C)+2)=O(G^{0.9})$. The terms of \eqref{eq:exponent} are bounded as follows:
\begin{align*}
  e\,U&\le4c_e\,G/\sqrt{\log G},\\
  (e+2)\,gP/K&\le3(4G+5g)/G^{1/300}=O(G^{1-1/300}),\\
  eK^2+C_0K^2\log(K(L+B'))&=O\bigl(G^{1/15}(\log G+\log(k(C)+2))\bigr)=O(G^{0.97}),\\
  \log\tau&\le2G/g+O(\log G+\log(k(C)+2))=O(G^{1-8/300}+G^{0.9}),\\
  g\log K&=O(G^{8/30}\log G),\\
  C_1\log(g(P+B))&=O(\log G+\log(k(C)+2))=O(G^{0.9}).
\end{align*}
Every term except the first is $o(G/\sqrt{\log G})$, uniformly in the circuit. The value $Z$ has
$O(U+gB)$ bits, so reducing and decoding it, and the preparation, cost $\poly(n+G+G_1)$. Hence the
total is $2^{cG/\sqrt{\log G}}\poly(n+G_1)$ for an absolute $c$ (any $c>4c_e$ works for large
$G$). This proves \cref{thm:main}.\qed

\begin{remark}[The single-qubit gate count]\label{rem:G1}
$G_1$ enters polynomially (input, merging) and through $k(C)\le G+G_1$, which sets the table bit
length $B=O(G/g+k(C))$; $B$ reaches the exponent only via $C_0K^2\log(K(L+B'))$, hence the
hypothesis $\log(G_1+2)\le G^{0.9}$. If it fails, the direct sum over $2G$ segment bits is
quasi-polynomial in the input length. The route through \A{} is exact whatever $k(C)$ is.
\end{remark}


% =====================================================================
\section{The route through \texorpdfstring{\Oa}{[O]}: proof of \texorpdfstring{\cref{thm:openai}}{Theorem 1.2}}\label{sec:openai}
% =====================================================================

We check the hypotheses of \cref{imp:tensor} one by one, with $s=\max\{4G,k(C),4\}$ and $C=4$.

\begin{itemize}
\item \emph{Bipartite, one-bit edges, at most $Cs$ edges.} The merged gates form one side. Every
segment is subdivided by a binary equality tensor, as \Oa{} does in its own CNF encoding; the
equality vertices form the other side. There are $2s_0\le4G\le s$ edges. Fixed legs ($x_q$,
$y_q$) are absorbed into the gate tensor, and a gate all of whose legs are fixed is an isolated
vertex.
\item \emph{Maximum degree at most $\lceil s^{9/10}\rceil$.} The degree is at most $4$, and
$\lceil s^{9/10}\rceil\ge4$ for $s\ge4$.
\item \emph{Number of vertices polynomial in $s$.} It is $G+s_0\le3G$.
\item \emph{Entries in $\cB_{P_0,\ell}$ with $1\le P_0,\ell\le Cs$.} Map $x\mapsto V$. The entries
of $\hat N_i$ become integer polynomials of degree at most $3$, reduced modulo $2^{P_0}$, and the
equality tensors have entries $0,1$. Over $\Z$ the contraction is $Q(V):=\sum_\sigma\prod_i\hat N_i(V)$, of
degree at most $3G$, so with $\ell:=3G+1\le s$ nothing is truncated and the contraction in
$\cB_{P_0,\ell}$ is $Q(V)\bmod 2^{P_0}$. Reducing $Q$ modulo $V^4+1$ over $\Z/2^{P_0}$ gives $A$
modulo $2^{P_0}$ coefficientwise, because $Q\bmod(x^4+1)=A$ in $\Zw$. By \cref{lem:galois},
$|a_j|\le2^{k_t/2}<2^{P_0-1}$ for $P_0:=\lceil k(C)/2\rceil+3$, so the balanced residues are the
$a_j$. Here $P_0\le s/2+4\le2s$. \emph{This is the only place where $k(C)$ enters, and it is why
$s\ge k(C)$}: \Oa{} works modulo a power of $2$ with precision at most $Cs$.
\item \emph{Entries computable in time $(s+\cL)^b2^{As^{9/10}}$ and space $(s+\cL)^b$.} Each entry is
a lookup in a merged gate, precomputed in polynomial time and space; $\cL=\poly(G+k(C))$.
\end{itemize}
The conclusion of \cref{imp:tensor} gives $(s+\cL)^{O(1)}2^{O(s/\sqrt{\log s})}$ bit operations and
$(s+\cL)^{O(1)}$ space; merging and the untouched scalars cost polynomial time and space.\qed

When $k(C)\gg G$ we see no way around the precision cap inside \Oa's ring: the numerator of a
two-qubit Clifford+T unitary can need arbitrarily many bits.

% =====================================================================
\section{Corollaries and the converse}\label{sec:corollaries}
% =====================================================================

\subsection{Marginals: the doubled network}
\begin{proof}[Proof of \cref{cor:marginals}, marginals]
Untouched lines contribute $|u_q|^2/2^{\kappa_q}$ for $q\in S$ and $1$ for $q\notin S$. On the
touched lines take two copies of the merged gates, the second with every entry conjugated
($\overline{a_0+a_1\omega+a_2\omega^2+a_3\omega^3}=a_0-a_3\omega-a_2\omega^2-a_1\omega^3$). Both
copies get input $x$. For $q\in S$ the output legs of both copies are fixed to $y_q$; for a touched
$q\notin S$ the two output legs are joined into one new \emph{trace} segment. The resulting sum over
paths has $2G$ merged gates and at most $2s_0+n_t\le4G$ segments, and its value is
$R=\sum_{y_{\bar S}}A(y)\overline{A(y)}=2^{k_t}p_t(y_S)$, where $p_t$ is the marginal of $C_t$.
$R$ is real, so $R=r_0+r_1\sqrt2$ with $r_0,r_1\in\Z$. For every $t$,
$\sigma_t(R)=2^{k_t}\sum_{y_{\bar S}}|\amp{y}{\sigma_t(C_t)}{x}|^2\in[0,2^{k_t}]$ (the proof of
\cref{lem:galois}), so $\|R\|_\infty\le2^{k_t}$ and $b:=k_t+2$ suffices in \cref{lem:kron}.
\Cref{sec:network,sec:time} now apply verbatim with $2G$ merged gates, $U\le8G$,
$P=4\lceil2G/g\rceil+1$ and $k_t$ replaced by $2k_t$; the exponent changes by a constant factor.
Under $(\mathrm{H}_{\mathrm{O}})$ take $s=\max\{8G,2k(C),4\}$, $\ell=6G+1$ and $P_0=k(C)+3$.
\end{proof}

\begin{proof}[Proof of \cref{cor:marginals}, sampling]
Draw $u\in[0,1)$ uniformly, revealing its binary digits only when needed, and output the
$y\in\{0,1\}^n$ whose interval in the lexicographic partition of $[0,1)$ into intervals of
lengths $|\amp{y}{C}{x}|^2$ contains $u$; then $\Pr[y]=|\amp{y}{C}{x}|^2$ exactly. The digits
$y_1,y_2,\dots$ are found one at a time: if the interval of the prefix $y_{<j}$ is
$[L,L+p(y_{<j}))$, then $y_j=0$ iff $u<L+p(y_{<j}0)$. This takes $n$ marginal computations,
with endpoints in $\Z[\sqrt2]/2^{k(C)}$, compared exactly with dyadic rationals (compare squares of
integers). A comparison is undecided after $t$ digits with probability at most $2^{-t}$, so
$\log n+O(1)$ digits are needed in expectation.
\end{proof}

\begin{proof}[Proof of \cref{cor:depth}]
$G\le dn/2$ and $G_1\le dn$. If $G\ge2$ and $\log(G_1+2)\le G^{0.9}$, use \cref{thm:main,cor:marginals}
and $G/\sqrt{\log G}\le dn/\sqrt{\log n}$ ($n\ge4$). Otherwise $G\le(\log(dn+2))^{1.12}$ and the
direct sum over at most $4G$ segment bits suffices.
\end{proof}

\subsection{The converse reduction}\label{sec:converse}
\begin{proof}[Proof of \cref{prop:tight}]
Make every variable occur at most three times, parsimoniously (Tovey's copying~\cite{Tovey84}): replace the $d\ge2$ occurrences of a
variable $v$ by fresh copies $v_1,\dots,v_d$ and add the $2$-clauses
$(\bar v_1\vee v_2),\dots,(\bar v_{d-1}\vee v_d),(\bar v_d\vee v_1)$, which force all copies to be
equal. The result $\varphi'$ has $n'=O(n+m)$ variables, $m'=O(n+m)$ clauses of width at most $3$,
each variable in at most three clauses, and $\#\varphi'=\#\varphi$.

Use one qubit per variable, one qubit $y_j$ per clause and one ancilla per clause, and put
$D:=\prod_jD_j$ with $D_j\ket{x,y}=(-1)^{y_j\cdot[\text{clause }j\text{ is violated by }x]}\ket{x,y}$.
Since $[\text{clause } j\text{ violated}]$ is the AND of the negated literals, $D_j$ is a
multi-controlled $Z$ on $y_j$ and at most three literal qubits, conjugated by $X$ on negative
literals. A controlled-controlled-controlled-$Z$ is exact over $\cG$ with one ancilla returned to
$\ket0$: compute the AND of two controls into the ancilla with a Toffoli gate, apply $\mathrm{CCZ}$,
and uncompute; Toffoli and $\mathrm{CCZ}$ have exact Clifford+T circuits with seven $T$ gates and
$O(1)$ $\mathrm{CNOT}$ and $H$ gates~\cite{NielsenChuang}. Every gadget touches $O(1)$ qubits and every variable
qubit lies in at most three gadgets, so the gadgets can be coloured greedily into $O(1)$ layers of
disjoint gadgets. Let $C_\varphi=H^{\otimes(n'+m')}DH^{\otimes(n'+m')}$, the Hadamards acting on
variable and clause qubits only. Then
\[
 \amp{0}{C_\varphi}{0}=2^{-(n'+m')}\sum_{x,y}\prod_j(-1)^{y_j[\text{clause }j\text{ violated}]}
 =2^{-(n'+m')}\sum_x\prod_j2[\text{clause }j\text{ satisfied}]=\frac{\#\varphi}{2^{n'}}.
\]
$C_\varphi$ has depth $O(1)$ and $O(n+m)$ gates. Reading off $\#\varphi$ needs only the amplitude.
\end{proof}

This construction is not scripted; the IQP analogue
$\amp{0}{H^{\otimes n}D_fH^{\otimes n}}{0}=\gap(f)/2^n$ for degree-$3$ $f$, with $\mathrm{CCZ}$
compiled into Clifford+T, is checked by \texttt{qc\_check.py} (\cref{sec:verification}).

% =====================================================================
\section{Fine-grained supremacy statements after the refutations}\label{sec:table}
% =====================================================================

Throughout this section ``refuted'' means \emph{refuted assuming the claimed refutations of ETH}.
Rows 1, 2, 5, 12 and 13 of \cref{tab:conj} need only the main theorem of either \A{} or \Oa,
which is what the labs have formalised in Lean.

\subsection{Two facts about SETH-type statements}
ETH fails if and only if $s_k=0$ for every $k\ge3$, where $s_k$ is the infimum of the exponents
$\delta$ with $k$-SAT in time $2^{\delta n}$~\cite{IPZ01}; then $s_\infty=\lim_ks_k=0$, so the
Strong Exponential Time Hypothesis $s_\infty=1$~\cite{ImpagliazzoPaturi01} fails too. Several
statements below are phrased for CNFs of unbounded width with at most $cn$ clauses. For the
decision version, Calabro, Impagliazzo and Paturi~\cite{CIP06} relate them to $k$-SAT through
Schuler's width reduction~\cite{Schuler05}. We need the counting version, which follows from the
same reduction because its two branches are disjoint.

\begin{lemma}[Counting width reduction]\label{lem:width}
Let $k\ge1$ and $s>0$, and suppose a deterministic algorithm counts the models of $k$-CNFs with $n'$
variables and $m'$ clauses in time $2^{sn'}\poly(n',m')$. Then the models of an arbitrary CNF with $n$
variables and $m$ clauses can be counted deterministically in time
$2^{sn}\exp\bigl(2^{-sk}(m+n/k)\bigr)\poly(n,m)$.
\end{lemma}

\begin{proof}
Write $\#_X F$ for the number of models of $F$ over the variable set $X$. If $F$ has a clause
$C=\{l_1,\dots,l_{k'}\}$ with $k'>k$, put $C'=\{l_1,\dots,l_k\}$, let $F_1=(F\setminus\{C\})\cup\{C'\}$, and
let $F_0$ be $F$ with $l_1,\dots,l_k$ set to false (satisfied clauses deleted, false literals
removed). Since $C'\subseteq C$, the models of $F_1$ are exactly the models of $F$ that satisfy $C'$;
the models of $F$ that violate $C'$ are the extensions of the models of $F_0$ by
$l_1=\dots=l_k=\text{false}$. These two sets are disjoint, so
$\#_XF=\#_XF_1+\#_{X\setminus\mathrm{var}(C')}F_0$. Recurse until no clause is longer than $k$ (a
leaf whose formula contains the empty clause counts $0$), and count each leaf with the given
algorithm. Along a root-to-leaf path every left branch removes a clause longer than $k$ and right
branches create none, so there are at most $m$ left branches, and a leaf reached by $r$ right
branches has $n-kr$ variables, so $r\le n/k$. The number of leaves reached by exactly $r$ right
branches is at most $\binom{m+r}{r}$~\cite[proof of Lemma~5]{CIP06}. Hence the time is at most
$\poly(n,m)\sum_{r\le n/k}\binom{m+r}{r}2^{s(n-kr)}\le\poly(n,m)\,2^{sn}\sum_{r}\binom{m+\lfloor n/k\rfloor}{r}2^{-skr}
\le\poly(n,m)\,2^{sn}(1+2^{-sk})^{m+n/k}$.
\end{proof}

\begin{corollary}\label{cor:density}
Assume the main theorem of \A{} or of \Oa. For all $c,\varepsilon>0$ there is a deterministic
algorithm that counts the models of CNFs (of any width) with $n$ variables and at most $cn$ clauses
in time $O(2^{\varepsilon n})$. Hence $\gap(f)=2^n-2\#f$ is computed in the same time.
\end{corollary}

\begin{proof}
Put $s=\varepsilon/3$ and choose $k$ with $2^{-sk}(c+1)\log e\le\varepsilon/3$. For this fixed
$k$, the main theorems count the models of $k$-CNFs in time $2^{o(n')}\poly(n',m')\le2^{sn'}\poly(n',m')$
for $n'$ large; \cref{lem:width} gives $2^{2\varepsilon n/3}\poly(n)$.
\end{proof}

\Cref{cor:density} does not give CNF-SAT with $m=n^{O(1)}$ clauses of unbounded width in time
$2^{(1-\delta)n}\poly(m)$: the term $2^{-sk}m$ then forces $k=\Omega(\log n)$, where neither main
theorem is subexponential. That form of SETH is not known to follow, and some statements below
rest on it.

\subsection{The table}
\Cref{tab:conj} lists the statements we found, with the wording checked against the primary
sources (arXiv versions cited). ``Witness family'' means the circuits built in the proof.
\emph{Refuted} means the statement itself is false under the refutations; \emph{premise refuted}
means a theorem whose hypothesis is now false, while its conclusion is not contradicted by
anything here; \emph{survives} means we know no consequence of the refutations that touches it.

{\footnotesize
\begin{longtable}{@{}>{\raggedright\arraybackslash}p{0.35cm}>{\raggedright\arraybackslash}p{7.3cm}>{\raggedright\arraybackslash}p{1.45cm}>{\raggedright\arraybackslash}p{6.6cm}@{}}
\caption{Fine-grained quantum-supremacy statements, assuming the refutations.}\label{tab:conj}\\
\toprule
& Statement (source) & Status & Reason\\
\midrule
\endfirsthead
\toprule
& Statement (source) & Status & Reason\\
\midrule
\endhead
\bottomrule
\endfoot
1 & \cite{MT19a} Conj.~1: ETH as $2^{\Omega(m)}$ for $3$-CNF decision, $m$ clauses & refuted & \A: $\#3$-SAT in $2^{O(m/\sqrt{\log m})}$.\\
2 & \cite{MT19a} Conj.~2 (``SETH''): deterministic CNF-SAT, $\le cn$ clauses, any width, needs $2^{(1-a)n}$ & refuted & \Cref{cor:density}.\\
3 & \cite{MT19a} Conj.~3 (NSETH, polynomial-size circuits) & survives & General circuits.\\
4 & \cite{MT19a} Conj.~4--6: nondeterministic $\gap\ne0$ for poly-size, log-depth, reversible circuits needs $2^{(1-a)n}$ & survive & $\poly(n)$-size instances; \cref{thm:main} gives nothing below $2^n$.\\
5 & \cite{MT19a} Conj.~7: nondeterministic $\gap(f)\ne0$, $3$-CNF with $m$ clauses, needs $2^{\Omega(m)}$ & refuted & $\gap(f)=2^n-2\#f$ in deterministic $2^{O(m/\sqrt{\log m})}$.\\
6 & \cite{MT19a} Thm.~6, \cite{HNS19} Thm.~1: ETH $\Rightarrow$ no $2^{o(t)}$ strong simulation of Clifford+T with $t$ $T$ gates (additive error $2^{-(3t+14)/7}$, resp.\ deciding $\amp0C0\ne0$) & premise refuted & Witness circuits have $O(t)$ gates and are solved exactly in $2^{O(t/\sqrt{\log t})}$; with $\poly(n)$ Cliffords the conclusion is open.\\
7 & \cite{MT19a} Thms.~7, 8: Conj.~7 $\Rightarrow$ no $2^{o(t)}$ multiplicative-error sampling of Clifford+T, resp.\ $t$-qubit PBC & premise refuted & Witness circuits ($O(t)$ gates) sampled exactly in $2^{O(t/\sqrt{\log t})}$; conclusions open.\\
8 & \cite{MT19a} Thm.~10: Conj.~2 $\Rightarrow$ exact probabilities with $h$ Hadamards not in $2^{(1-a)h}$ & premise refuted & Witness circuits have $\Theta(ch^2)$ gates; conclusion open.\\
9 & \cite{MT19a} Thm.~12: ETH $\Rightarrow$ efficient stabilizer rank of $t$-Toffoli resources is $2^{\Omega(t)}$ & premise refuted & Conclusion open.\\
10 & \cite{HNS18} Thm.~3: CNF-SAT needs $2^{n-o(n)}\poly(m)$ $\Rightarrow$ strong simulation of $\poly(n)$-size circuits to precision $2^{-n}/2$ needs $2^{n-o(n)}$ & survives & Premise is the CNF-SAT form with $\poly(m)$ clauses, not known to follow; unconditional monotone bound unaffected.\\
11 & \cite{HNS19} Thm.~2: an $O(2^{2.2451\cdot10^{-8}N})$ simulator ($N$ $T$ gates) gives $3$-SAT in $O(1.3^n)$ & moot & The refutations give the consequent outright.\\
12 & \cite{MT19b} Conj.~1 (``SETH'', $\le cn$ clauses, deterministic decision) & refuted & \Cref{cor:density}; Thm.~1 loses its premise.\\
13 & \cite{MT19b} Conj.~2: nondeterministic $\gap(f)\ne0$, CNF with $\le cn$ clauses, needs $2^{(1-a)n}$ & refuted & \Cref{cor:density} (counting form). Thm.~2 loses its premise; its witness circuits have $\ge nm$ qubits.\\
14 & \cite{MT19b} Conj.~3--4 (OV), 5--6 ($3$-SUM) & survive & OV is implied by SETH, not equivalent; $3$-SUM unrelated.\\
15 & \cite{DHKL20} poly3-NSETH$(a)$ & survives & Worst case has $\Theta(n^3)$ monomials (\cref{sec:cubic}).\\
16 & \cite{DHKL20} per-int-NSETH$(b)$ & survives & Dense permanents.\\
17 & \cite{DHKL20} poly3-ave-SBSETH$(a')$ & survives & Average case, dense, additive error.\\
18 & \cite{DHKL20} Thm.~2: ETH $\Rightarrow$ deterministic poly3-NONBALANCED needs $2^{an}$ & premise refuted & Reduction instances have $O(n)$ monomials; worst case open.\\
19 & \cite{MT20} Conj.~1: random cubic, multiplicative error, $\mathrm{NTIME}[n^2]$ oracle, $2^{an}$ & survives & Dense, average case, approximate.\\
20 & \cite{MT20} Conj.~2: random quadratic plus worst-case log-depth circuit & survives & General log-depth circuits.\\
21 & \cite{MT20} Conj.~3: random quadratic plus worst-case $3$-CNF with $m$ clauses, $2^{am}$ & survives (not shown safe) & $\Theta(n^2+m)$ gates; nothing for $m=O(n)$.\\
22 & \cite{AC17} QUATH (poly time, random circuits with $\ge n^2$ gates) and HOG & unaffected & Average case, approximate.\\
23 & \cite{BPS19} basic QSETH: no quantum $O(2^{n(1-\delta)/2}m^{O(1)})$ CNF-SAT & survives & CNF-SAT form; the $k$-CNF variant of their footnote~8 is refuted.\\
24 & $2^{\Omega(t)}$ $T$-count hardness with unboundedly many Cliffords; $\chi(\ket{T}^{\otimes t})=2^{\Omega(t)}$~\cite{BSS16} & open & ETH support gone; best known $2^{0.3963t}\poly$~\cite{QPG21}; compressed circuits have $\Theta(t^2)$ gates.\\
\end{longtable}
}

\emph{Row 13} was left open in our working notes: there we checked only the decision form of the
width--density duality. \Cref{lem:width} settles it.

\subsection{Sampling versus exact amplitudes}
By \cref{cor:marginals}, any hardness claim for sampling, with multiplicative or additive error,
stated as $2^{\Omega(n)}$ for a family with $G=o(n\sqrt{\log n})$ gates is refuted, since exact
samples are approximate ones. Nothing is said about families with $G\ge n\sqrt{\log n}$, which
include every random-circuit sampling proposal at depth $d\gg\sqrt{\log n}$, nor about
polynomial or unconditional separations for shallow circuits.

The estimates of~\cite{DHKL20} ($208$ IQP qubits and $500$ gates, assuming poly3-NSETH$(1/2)$ and a
cost linear in the number of gates) extrapolate from the dense worst case to a sparse family, which
is asymptotically $2^{O(n/\sqrt{\log n})}$ under the refutations; at $n=208$ this has no force.

% =====================================================================
\section{Cubic phases}\label{sec:cubic}
% =====================================================================

Let $f=\sum_{t\in T}m_t$ over $\mathbb F_2$, with $m_t$ monomials of degree at most $3$ in $n$
variables and $|T|=m$, and $\gap(f)=\sum_x(-1)^{f(x)}$. This is the object of poly3-NSETH
in~\cite{DHKL20}, and $\amp{0}{H^{\otimes n}D_fH^{\otimes n}}{0}=\gap(f)/2^n$ for the diagonal
circuit $D_f$ built from $Z$, $\mathrm{CZ}$ and $\mathrm{CCZ}$. Compiling $\mathrm{CCZ}$ into
Clifford+T gives $G=\Theta(m)$, so \cref{thm:main} computes $\gap(f)$ in $2^{O(m/\sqrt{\log m})}$,
which is $2^{o(n)}$ only for $m=o(n\sqrt{\log n})$. For CNFs, \A{} has a second bound that does
not depend on $m$. Does one hold for cubic phases? As far as we can see, no.

\subsection{Why the clause-count-free mechanism does not transfer}
As we read \A[Appendix~D], the $m$-free bound writes $\Ind{\text{no clause is violated at }\xi}$
by inclusion--exclusion over the set $V(\xi)$ of violated clauses, truncates it at a number of
marks of order $n/\log\Delta$, and handles the assignments with more violated clauses separately,
using that they contribute \emph{zero} to $\#\mathrm{SAT}$. For phases there is no such set:
$(-1)^{f(x)}=\prod_t(1-2m_t(x))=\sum_{\Pi\subseteq V(x)}(-2)^{|\Pi|}=(-1)^{|V(x)|}$, where $V(x)$ is
the set of monomials equal to $1$ at $x$, so every assignment contributes $\pm1$, and the
assignments with many ``violated'' monomials carry a non-zero total in general. This is checked
on random instances (check C4 in \cref{sec:verification}). What does transfer is a weaker,
$2$-adic truncation.

\begin{proposition}[$2$-adic truncation; proved]\label{prop:2adic}
For $S\subseteq T$ let $V(S)$ be the set of variables occurring in the monomials of $S$.
\begin{enumerate}[label=\textup{(\alph*)}]
\item $\gap(f)=\sum_{S\subseteq T}(-2)^{|S|}\,2^{n-|V(S)|}$.
\item Every term with $|S|\ge|V(S)|+2$, in particular every term with $|S|\ge n+2$, is divisible by
$2^{n+2}$. Since $|\gap(f)|\le2^n$, $\gap(f)$ is the representative in $(-2^{n+1},2^{n+1}]$ of
$\sum_{|S|\le n+1}(-2)^{|S|}2^{n-|V(S)|}$ modulo $2^{n+2}$.
\end{enumerate}
\end{proposition}

\begin{proof}
(a) Expand $\prod_t(1-2m_t(x))$ and use $\sum_x\prod_{t\in S}m_t(x)=2^{n-|V(S)|}$. (b) The $2$-adic
valuation of the term of $S$ is $|S|+n-|V(S)|$.
\end{proof}

The truncation at $n+1$ marks is sharp up to a constant factor: truncating at $\lfloor n/2\rfloor$
marks gives wrong answers on some instances (check C3). So the budget is $\Theta(n)$ marks, where
\A{} pays about $n/\log\Delta$.

\emph{Sketch (not proved here).} Feeding (b) into the incidence encoding that \Oa{} uses for sparse
$\mathbb F_2$-polynomial systems (each monomial a test vertex of weight $1-2\prod_{v\in t}z_{t,v}$,
precision $2^{n+2}$, $\ell=1$, incidence degree bound $3(n+1)$, and branching on the variables in
more than $(3m/n)\sqrt{\log n}$ monomials) should give $O(n\log(em/n)+\sqrt n\log n)$ wires, hence,
under $(\mathrm{H}_{\mathrm{O}})$, $\gap(f)$ in time $2^{O(n\log(em/n)/\sqrt{\log n})}\poly(m)$:
modulo $2^{n+2}$ only sets of at most $n+1$ marked monomials matter, at about $\log(em/n)$ bits
each. We have not checked this application line by line against the
degree and entry hypotheses of \Oa{} (the maximum degree should be $n^{3/4}2^{o(\sqrt{\log n})}\le s^{9/10}$
and the entries are products of $1$ and $-2$, so we expect no obstacle). If correct, it is $2^{o(n)}$ if and
only if $m\le n2^{o(\sqrt{\log n})}$, more than \cref{thm:main} gives, and it is worse than brute
force for dense $f$ ($m=\Theta(n^3)$ gives $2^{O(n\sqrt{\log n})}$). Making cubic phases
$m$-independent would decide dense poly3-NONBALANCED in $2^{o(n)}$ and refute poly3-NSETH$(a)$
for every $a$; we record it as a target, not a result.

% =====================================================================
\section{Limitations}\label{sec:limitations}
% =====================================================================

\begin{enumerate}[label=\textup{(\alph*)}]
\item \emph{Everything algorithmic is conditional} on $(\mathrm{H}_{\mathrm{A}})$,
$(\mathrm{H}_{\mathrm{O}})$, or (for \cref{cor:density}) the main theorem of either manuscript.
Expert review of the manuscripts is pending; if the hypotheses fail, the results lose their proofs,
and we offer no independent route. The imported statements are used in sharp forms that are not
machine-checked in the generality needed here; our proofs are on paper, written with AI
assistance, and not yet checked by a human.
\item \emph{Exact, gate-count-limited, galactic.} The results concern exact simulation of
circuits with $G=o(n\sqrt{\log n})$ gates; approximate and average-case hardness (random-circuit
sampling, cross-entropy benchmarking, QUATH) and polynomial-depth families are untouched, and the
crossover with the state vector is astronomically large (\cref{sec:intro}).
\item \emph{Denominators in the route through \Oa.} Its $2$-adic precision cap forces $k(C)=O(s)$.
\item \emph{Sketches.} The bound after \cref{prop:2adic} is a sketch; the circuit of
\cref{prop:tight} is not scripted; the reading of \A[Appendix~D] is ours.
\item \emph{Literature.} Every row of \cref{tab:conj} was checked against the arXiv text of its
source; publication venues are given only where we could confirm them. We have not searched
systematically for prior statements of \cref{thm:main}; given how direct it is, it may well have
been noticed independently.
\item \emph{Scale of the numerical checks.} They cover $n\le6$ and $G\le20$ and test the reductions
(encodings, decodings, parameters, ring emulation), not the engines, which are replaced by brute
force and variable elimination.
\end{enumerate}

% =====================================================================
\section{Verification}\label{sec:verification}
% =====================================================================

\subsection{Numerical checks}
The scripts (\texttt{research/scripts/qc\_*.py} in the project repository) use exact integer and
$\Zw$ arithmetic only and print PASS or FAIL; \texttt{qc\_lib.py} implements the objects of
\cref{sec:proof,sec:openai,sec:corollaries}.

\texttt{qc\_check.py [seed] [count]} draws random circuits ($n\in[2,6]$, $G\in[1,20]$,
$G_1\in[0,25]$, random gates, $x$, $y$, and balanced or contiguous partitions into $g\le4$ groups)
and checks: (Q1) \cref{lem:paths} against the state vector; (Q2) the integer network by brute
force over the crossing bits, by the literal sum of \cref{def:network} on small cases and by
variable elimination, decoded by \cref{lem:kron}; (Q3) decoding of the unreduced Kronecker
polynomial; (Q4) condition~(i), $U\le4G$ and $|W_a|<2^B$; (Q5) the contraction in
$\cB_{P_0,\ell}$ with $P_0=\lceil k/2\rceil+3$, $\ell=3G+1$, reduced modulo $V^4+1$ and lifted;
(Q6) \cref{lem:galois}, with $|\sigma_t(A)|^2\le2^k$ checked exactly in $\Z[\sqrt2]$; (Q7) marginals
via the doubled network, in $\Z[\sqrt2]$ and equal to $1$ for $S=\emptyset$. It also checks $20$
IQP instances ($\gap(f)/2^n$ with $\mathrm{CCZ}$ compiled into Clifford+T) and four ablations that
must be caught: $b$ one bit too small, $\ell=4$, $P_0$ one bit short, and no conjugation in the
doubled network. Seed~$1$ ($150$ circuits): $529$ networks by brute force, $418$ also by the literal
sum, $65$ truncated-ring brute forces, $91$ marginals, at most $23$ crossing segments, $48$
partitions with a node of more than $14$ boundary bits skipped. Seed~$2$: $544$, $432$, $62$, $86$,
$24$, $39$. Both PASS, with all ablations caught; the agreement is equality of integers.

\texttt{qc\_cubic\_gap.py} checks, for random $f$ with $n\le7$, $m\le14$: (C1)--(C2)
\cref{prop:2adic}; (C3) the valuation bound, and failure of truncation at $\lfloor n/2\rfloor$ marks
on some instance; (C4) a non-zero total of the assignments with more than $\lfloor n/3\rfloor$
``violated'' monomials on some instance. Seeds $1$ ($120$ instances) and $3$ ($400$) PASS.

\subsection{Formalisation}\label{sec:formalisation}
\Cref{thm:main} is formally verified in Lean~4 on top of \A's development: \texttt{amplitude\_TM\_rate'}
(the stated bound under the condition on $L$), \texttt{amplitude\_TM\_rate} (all circuits) and
\texttt{amplitude\_TM} (every $\delta>0$), in the public repository~\cite{Tulsyan26} (directories
\texttt{lean/EthFollowup/Quantum} and \texttt{lean-bridge-anthropic/bridge/EthFollowup/Quantum}).
The semantics a reader must check are short: the gate matrices over $\mathbb{C}$, the circuit matrix as
the ordered product, and the amplitude $\amp{y}{C}{x}$ (\texttt{Circuit.lean}, \S1), and the input
and output encodings (\S7). The theorems depend only on the axioms \texttt{propext},
\texttt{Classical.choice} and \texttt{Quot.sound}, and fresh \texttt{leanchecker} replays of their
closures (\texttt{EthFollowup.Quantum.QRate}, which contains \texttt{QTheorem}, and
\texttt{EthFollowup.Quantum.Paths}) passed. \Cref{thm:openai}, \cref{cor:marginals} and \cref{cor:density} are proved on paper
only and are not claimed as machine-checked.

\section*{Acknowledgements}
This work was done with AI assistance (Claude Code, Anthropic): the observation, the proofs, the
literature checks and the numerical verification scripts were produced in sessions with Claude
Code. Everything here has been checked numerically or against the primary sources as described,
but awaits human expert review. Responsibility for any errors rests with the author.

\bibliographystyle{amsplain}
\bibliography{refs}

\end{document}
